\documentclass[11pt,a4paper]{article}

\usepackage[T2A]{fontenc}
\usepackage[utf8]{inputenc}
\usepackage[serbian]{babel}

\usepackage[a4paper,margin=2.5cm]{geometry}

\usepackage{amsmath,amssymb,amsthm}
\usepackage{tikz}
\usepackage{tasks}
\usepackage{enumitem}
\usepackage{hyperref}
\usepackage{xcolor}
\usepackage{tcolorbox}

\usetikzlibrary{calc}

% ---------------- BOJE ----------------

\definecolor{zbg}{RGB}{235,255,235}
\definecolor{zfr}{RGB}{20,120,60}

\definecolor{tbg}{RGB}{245,235,255}
\definecolor{tfr}{RGB}{110,60,170}

% ---------------- BOKSEVI ----------------

\newtcolorbox{zadatakBox}{
    colback=zbg,
    colframe=zfr,
    title=Задатак
}

\newtcolorbox{theoremBox}{
    colback=tbg,
    colframe=tfr,
    title=Теорема
}

% ---------------- TEOREME ----------------

\newtheorem{theorem}{Теорема}

\begin{document}

\begin{center}
{\LARGE \textbf{Питагорина теорема}}
\end{center}

\begin{theoremBox}
\begin{theorem}[Питагорина теорема]
\label{thm:pitagora}
За правоугли троугао важи:
\[
a^2 + b^2 = c^2
\]
где су $a$ и $b$ катете, а $c$ хипотенуза.
\end{theorem}
\end{theoremBox}

\textbf{Доказ:}

Посматрајмо квадрат странице дужине $a+b$. Идеја је да приметимо да се унутар квадрата налазе четири међусобно подударна правоугла троугла и централни квадрат странице дужине $c$.
\begin{center} 
\begin{tikzpicture}[scale=0.8]

\draw[thick] (0,0) rectangle (8,8);

\coordinate (A) at (0,5); 
\coordinate (B) at (5,8); 
\coordinate (C) at (8,3); 
\coordinate (D) at (3,0); 

\draw[thick] (A)--(B)--(C)--(D)--cycle;

\draw[thick] (0,0)--(A);
\draw[thick] (0,8)--(B); 
\draw[thick] (8,8)--(C); 
\draw[thick] (8,0)--(D); 

\node at (2.2,2.2) {$c$}; 
\node[left] at (0,2.5) {$a$}; 
\node[left] at (0,6.5) {$b$}; 
\node[below] at (2,0) {$b$};
\node[below] at (6.5,0) {$a$}; 

\end{tikzpicture} 
\end{center} 

Површина великог квадрата једнака је: \[ (a+b)^2 \] 

Са друге стране, површина тог квадрата једнака је збиру:

\begin{itemize} 
\item површина четири правоугла троугла,
\item површине унутрашњег квадрата. 
\end{itemize} 

Дакле:
\[ (a+b)^2=4\cdot \frac{ab}{2}+c^2 \]

Сређивањем добијамо: \[ a^2+2ab+b^2=2ab+c^2 \] 
\[ a^2+b^2=c^2 \]

Чиме је теорема доказана. \hfill $\square$

\vspace{0.5cm}

\begin{theoremBox}
\begin{theorem}[Обратна теорема]
Нека је дат троугао са страницама $a,b$ и $c$. Ако важи $a^2 + b^2 = c^2$, троугао је правоугли.
\end{theorem}
\end{theoremBox}

\textbf{Доказ.} 

Нека је дат троугао са страницама $a$, $b$ и $c$ тако да важи \[ a^2+b^2=c^2. \] 

Конструишемо правоугли троугао чије су катете дужине $a$ и $b$ (њих лако препознамо као две мање странице).
Према Питагориној теореми~\ref{thm:pitagora}, дужина његове хипотенузе износи \[ \sqrt{a^2+b^2}. \] 

Пошто је \[ a^2+b^2=c^2, \] следи \[ \sqrt{a^2+b^2}=\sqrt{c^2}=c. \] 
Дакле, конструисани правоугли троугао заправо има странице $a$, $b$ и $c$. 

Према ставу подударности троуглова ССС, тај троугао је подударан са полазним (датим) троуглом. 

Како је конструисани троугао правоугли, следи да је и полазни троугао правоугли. 

\hfill $\square$

\vspace{0.5cm}

% ---------------- ZADATAK ----------------

\begin{zadatakBox}
Једнакокраки троугао има висину дужине 12 и основицу дужине 10. Наћи дужину крака једнакокраког троугла.
\end{zadatakBox}

\textbf{Решење:}
Висина дели једнакокраки троугао на два подударна правоугла троугла. Применом Питагорине теореме добијамо
\[
h_o^2+(\frac{o}2)^2=c^2  \iff  12^2+5^2=c^2 \iff   144+25=c^2 \iff   c^2 = 169 \iff  c = 13.
\]
што нам даје дужину оба крака. Решење се лако чита и са следеће слике: 




\begin{center} 
\begin{tikzpicture}[scale=0.8]

%\draw[thick] (0,0) rectangle (8,8);

\coordinate (A) at (0,0); 
\coordinate (B) at (8,0); 
\coordinate (C) at (4,8); 
\coordinate (D) at (4,0); 

\draw[thick] (A)--(B)--(C)--(D)--(A)--(C);

\node[left] at (A) {$A$};
\node[right] at (B) {$B$}; 
\node[above left] at (C) {$C$};

\node[left] at ($(A)!0.5!(C)$) {$c$};
\node[right] at ($(B)!0.5!(C)$) {$c$}; 
\node[below right] at ($(A)!0.5!(B)$) {$o$};
\node[right] at ($(C)!0.7!(D)$) {$h_o = 12$};

\node[below] at ($(A)!0.5!(D)$) {$5$};
\node[below] at ($(B)!0.5!(D)$) {$5$};

\draw (5,1) rectangle (4,0);



\end{tikzpicture} 
\end{center}


\vspace{0.5cm}


% ---------------- TABELA PITAGORINIH TROJKI ----------------

\section*{Питагорине тројке} 

\begin{center} 
\renewcommand{\arraystretch}{1.4} 

\begin{tabular}{|c|c|c|} 
\hline \textbf{Катета} & \textbf{Катета} & \textbf{Хипотенуза} \\ 
\hline 3 & 4 & 5 \\ 
\hline 5 & 12 & 13 \\ 
\hline 8 & 15 & 17 \\ 
\hline 7 & 24 & 25 \\ 
\hline $m^2-n^2$ & $2mn$ & $m^2+n^2$ \\
\hline \end{tabular} \end{center} 

Последња врста представља општи облик за генерисање Питагориних тројки за два произвољна природна броја.
\footnote{јер важи $(m^2 - n^2)^2 + (2mn)^2 = m^4 - 2m^2n^2 + n^4 + 4m^2n^2 = m^4 + 2m^2n^2 + n^4 = (m^2 + n^2)^2$}

% ---------------- 3 TEST PITANJA ----------------
%\newpage
\section*{Тест}

\begin{enumerate}[label=\arabic*.]

\item Тврђење Питагорине теореме је:
\begin{tasks}(2)
\task $a+b=c$
\task $a^2+b^2=c^2$
\task збир квадрата над катетама једнак је квадрату над хипотенузом
\task $(a  + b)^2=c^2$
\end{tasks}

\item За катете дужине 3 и 4, хипотенуза је дужине:
\begin{tasks}(4)
\task $5$
\task $6$
\task $8$
\task $10$
\end{tasks}

\item Које од наведених тројки су Питагорине:
\begin{tasks}(4)
\task $(3, 4, 5)$
\task $(8, 15, 16)$
\task $(8, 8, 8)$
\task $(12, 13, 5)$
\end{tasks}

\end{enumerate}

\end{document}